\section{Однооборотное отображение}
\label{sec:map}

Главы~\ref{sec:magnetic}--\ref{sec:electrostatic} получают $\delta_{\mathrm{eq}}$ усреднением периода по синхротронной фазе и по бетатронной фазе. То же число есть среднее $\delta$ вдоль орбиты повторения однооборотного отображения с высокочастотным зазором. Ниже это отображение записывается в координатах COSY~Infinity~\cite{Berz1990}, и из него собираются те же $\eta_0$, $\eta_1$ и $\Lambda_T$, что уже стоят в~\eqref{eq:beta-deq} и~\eqref{eq:es-deq}. Формулы предыдущих глав не меняются.

\subsection{Равновесие как среднее по орбите отображения}
\label{sec:map-orbit}

Без зазора шестая координата постоянна, и отображение одного оборота переводит поперечные координаты и время. Зазор делает отображение шестимерным: за один проход он сдвигает энергию на $eV\sin\varphi$, а фаза зазора видит опоздание $T-T_s$. Обозначим это отображение $M$. Захваченная частица есть орбита повторения $M$, то есть последовательность точек $z_{n+1}=M(z_n)$, лежащая на инвариантном торе трёх действий. Равновесный импульс этой главы есть среднее шестой координаты по этому тору. Пока амплитуда синхротронных колебаний мала, среднее совпадает с положением центра тора по переменной, сопряжённой времени.

Линейная часть $M$ этого центра не содержит. Сдвиг $\delta_{\mathrm{eq}}$ квадратичен по амплитудам: он рождается членами $\delta^2$ и $J$ в периоде и становится виден только во втором порядке нормальной формы. Усреднение глав~\ref{sec:magnetic}--\ref{sec:electrostatic} есть вычисление этого центра. Матрица одного оборота задаёт тюны и Твисс и к $\delta_{\mathrm{eq}}$ не приводит.

Координаты COSY выбраны так, что отображение симплектично~\cite{Berz1990}. Пятая координата есть
\begin{equation}
\label{eq:map-l}
l=-(t-t_0)\,v_0\,\frac{\gamma_s}{\gamma_s+1},
\end{equation}
шестая есть $\delta_K=(K-K_0)/K_0$, относительная кинетическая энергия в опорной точке. На целом обороте отсюда следует уже записанное~\eqref{eq:es-cosy-T}. Высокочастотное поле видит $T$, а не путь: в период входит $l$, и магнитная пересчётная формула к $l$ второй раз не применяется.

\subsection{Коэффициенты отображения и $\eta$, $\Lambda_T$}
\label{sec:map-coef}

Без зазора $\delta_K$ постоянно, и время одного оборота есть пятая компонента отображения. Коэффициенты Тейлора COSY суть множители мономов: оператор извлечения уже делит производную на факториалы показателей, так что $l\supset(5|6)\,\delta_K+(5|66)\,\delta_K^2$. Линейные множители совпадают с элементами матрицы, $(5|6)=\mathrm{ME}(5,6)$.

Этот ряд, взятый при $x=a=0$, не есть время замкнутой орбиты. Частица, выпущенная из начала координат с $\delta_K\neq 0$, колеблется вокруг дисперсионной орбиты $x=D\delta_K$. Время оборота, которое входит в $\eta_0$ и $\eta_1$, есть $l$ на самой замкнутой орбите. Пусть при фиксированном $\delta_K$ её начало в точке наблюдения есть $x_*=D_1\delta_K+D_2\delta_K^2$ и $a_*=D_1'\delta_K+D_2'\delta_K^2$. Подстановка в многочлен пятой компоненты даёт мономы замкнутой орбиты
\begin{align}
e_1
&=
(5|6)+(5|1)D_1+(5|2)D_1',
\label{eq:map-e1}\\
e_2
&=
(5|66)+(5|1)D_2+(5|2)D_2'
+(5|16)D_1+(5|26)D_1'
\nonumber\\
&\quad
+(5|11)D_1^2+(5|12)D_1 D_1'+(5|22)(D_1')^2.
\label{eq:map-e2}
\end{align}
Процедура \texttt{ECM} делает ту же подстановку: она ищет поперечную неподвижную точку и возвращает $l_{\mathrm{out}}-l_{\mathrm{in}}$ как функцию канонического $\delta_K$. Коэффициенты $e_1$ и $e_2$ суть её мономы по $\delta_K$. Именно их надо подставлять в~\eqref{eq:es-eta0-map} и~\eqref{eq:es-eta1-map} на место $(5|6)$ и $(5|66)$. Вместе с рядом~\eqref{eq:es-dk-series} они дают
\begin{align}
\eta_0
&=
-\left(\frac{\gamma_s+1}{\gamma_s}\right)^2\frac{e_1}{C_s},
\label{eq:map-eta0}\\
\eta_1
&=
-\frac{\gamma_s+1}{\gamma_s C_s}
\Biggl(
e_1\,\frac{\gamma_s+1}{2\gamma_s^3}
+e_2\left(\frac{\gamma_s+1}{\gamma_s}\right)^2
\Biggr).
\label{eq:map-eta1}
\end{align}
Это $\eta_0$ и $\eta_1$ периода, магнитные или электростатические, смотря по полю, которое породило отображение. На игрушечном симплектическом кольце из тонких гамильтоновых киков производная времени по $\delta$ в начале координат и на замкнутой орбите различаются на $0{,}21$ относительно значения на орбите: поправка~\eqref{eq:map-e1} не является малой высшего порядка.

Бетатронное слагаемое собирается из тех же мономов пятой компоненты. На линейном эллипсе Твисса в точке наблюдения, $\langle x^2\rangle=J\beta$, $\langle xa\rangle=-J\alpha$, $\langle a^2\rangle=J\gamma_T$, среднее квадратичной формы есть
\begin{equation}
\label{eq:map-ell}
\frac{\langle\Delta l\rangle_{\mathrm{ell}}}{J}
=
(5|11)\,\beta-(5|12)\,\alpha+(5|22)\,\gamma_T
\end{equation}
и так же по вертикали. Один коэффициент $(5|11)$ средним по фазе не является: в него не входят $\langle xa\rangle$ и $\langle a^2\rangle$. Форма~\eqref{eq:map-ell} есть среднее по эллипсу, описанному вокруг начала координат. Инвариантная кривая нелинейного отображения сдвинута. Квадратичный толчок поперечного отображения, усреднённый по тому же эллипсу, даёт за оборот приращение $(\langle\Delta x\rangle,\langle\Delta a\rangle)$, пропорциональное $J$. Центр, который это приращение уравновешивает, есть
\begin{equation}
\label{eq:map-centre}
\begin{pmatrix}x_c\\ a_c\end{pmatrix}
=
(I-R)^{-1}
\begin{pmatrix}\langle\Delta x\rangle\\ \langle\Delta a\rangle\end{pmatrix},
\end{equation}
где $R$ --- поперечная матрица одного оборота. Линейная связь времени с координатой добавляет $(5|1)x_c+(5|2)a_c$. Вертикальное действие тем же способом сдвигает и вертикальный, и, через члены $y^2$ горизонтального толчка, горизонтальный центр; в приращение $l$ входят все четыре линейных коэффициента $(5|1),\ldots,(5|4)$. Сумма
\begin{equation}
\label{eq:map-LT}
\frac{\partial\langle\Delta l\rangle}{\partial J}
=
\frac{\langle\Delta l\rangle_{\mathrm{ell}}}{J}
+(5|1)\frac{x_c}{J}+(5|2)\frac{a_c}{J}
\end{equation}
не зависит от точки наблюдения. По отдельности квадратичная форма от точки зависит: на игрушечном кольце при одном и том же секступоле отношение формы к $2\pi\partial\nu/\partial\delta$ равно $0{,}853$ в фокусе и $0{,}881$ в дефокусе, а сумма~\eqref{eq:map-LT} в обеих точках равна $2\pi\partial\nu/\partial\delta$ с остатком $10^{-5}$. Переход к периоду снова есть~\eqref{eq:es-cosy-T}, и $\Lambda_T$ получается умножением~\eqref{eq:map-LT} на $-(\gamma_s+1)/(\gamma_s C_s)$. Это уточнение~\eqref{eq:es-map-L}: там стоит эллипс, здесь к нему добавлен центр. На кольце без квадратичных поперечных членов центр равен нулю, и две записи совпадают.

Та же сумма есть смешанная производная нормальной формы. Пусть $K(J,\delta)$ --- гамильтониан, чьему сдвигу на один оборот отвечает симплектическое отображение в переменных, где $\mathrm{d}\tau=\partial K/\partial\delta$ и $\mathrm{d}\theta=\partial K/\partial J$, а $J$ --- действие с площадью эллипса $2\pi J$. Тогда $\partial\langle\Delta\tau\rangle/\partial J=\partial\mu/\partial\delta=2\pi\,\partial\nu/\partial\delta$. В координатах COSY $\tau=l$ и $\delta=\delta_K$, поэтому
\begin{equation}
\label{eq:map-id}
\frac{\partial\langle\Delta l\rangle}{\partial J}
=
2\pi\frac{\partial\nu}{\partial\delta_K}.
\end{equation}
Хроматичность глав~\ref{sec:magnetic}--\ref{sec:electrostatic} есть $\xi=\partial\nu/\partial\delta$ по импульсу на нулевом потенциале. На нуле $\mathrm{d}\delta_K/\mathrm{d}\delta=(\gamma_s+1)/\gamma_s$. Множитель $(\gamma_s+1)/\gamma_s$ в~\eqref{eq:es-cosy-T} и в этой производной сокращается, и
\begin{equation}
\label{eq:map-xi}
\frac{\Lambda_T}{J}
=
-\frac{2\pi\xi}{C_s}.
\end{equation}
Это тождество~\eqref{eq:es-sym} для той пары переменных, в которой оно выведено. Оно относится к среднему по инвариантной кривой, то есть к сумме~\eqref{eq:map-LT}, а не к одному эллипсу~\eqref{eq:map-ell}.

На игрушечном кольце тождество проверяется прямо. Пока секступоль выключен, центр равен нулю и уже одна форма~\eqref{eq:map-ell} воспроизводит $2\pi\partial\nu/\partial\delta$. Скан силы секступоля по одной квадратичной форме даёт наклон $1{,}296$ и ненулевой свободный член. Та же прямая по сумме~\eqref{eq:map-LT} имеет наклон $0{,}99993$ и свободный член, меньший $10^{-3}$ при наклонах порядка $10$. Если у квадрупольного кика отбросить продольную половину, якобиан на смещённой точке перестаёт быть симплектическим, отношение при нулевом секступоле падает до $1/2$, и свободный член прямой остаётся большим; наклон по секступолю центр по-прежнему возвращает к единице, потому что собственный вклад секступоля в $l$ кубичен по координате и в порядок $J$ не входит. Продольная половина квадрупольного кика сидит в квадратичной форме и центром не заменяется.

Отсюда читается коэффициент $\kappa$ главы~\ref{sec:electrostatic}. Там $\Lambda/J$ измерялось одним оборотом по линейному эллипсу, без~\eqref{eq:map-centre}. На текущем тексте модели `_verify_ch3_ring.py` такая прямая по магнитному скану имеет $\kappa=1{,}065$ и $b=5{,}2\times 10^{-4}\,\mathrm{м}^{-1}$; после добавления центра $\kappa=0{,}990$ и $b=-4{,}3\times 10^{-4}\,\mathrm{м}^{-1}$. Это тот же приём, который в тексте главы дал $\kappa=1{,}076$: расхождение с единицей есть пропущенный центр, а не множитель в $\delta_{\mathrm{eq}}$. Естественная хроматичность этого запуска равна $-0{,}138$, а в главе записано $-0{,}207$; таблица главы поэтому не пересчитывалась, и сюда переносится механизм, а не новая подгонка взамен той таблицы. На цилиндре того же скрипта эллипс даёт $\kappa=1{,}154$ и $b=-8{,}8\times 10^{-3}\,\mathrm{м}^{-1}$, центр приводит к $\kappa=0{,}940$ и $b=-2{,}2\times 10^{-3}\,\mathrm{м}^{-1}$. При секступоле, где $\xi=0{,}006$, эллипс ещё держит наклон $-5{,}8\times 10^{-3}\,\mathrm{м}^{-1}$, а сумма с центром равна $+3\times 10^{-4}\,\mathrm{м}^{-1}$ рядом с $-2\pi\xi/C_s=-2\times 10^{-4}\,\mathrm{м}^{-1}$. Большой остаток при погашенной хроматичности в основном был пропущенным центром. То, что на цилиндре после центра остаётся $\kappa=0{,}94$, есть свойство жёсткого трекера со скачком потенциала: игрушечное симплектическое отображение такого остатка не имеет. В $\delta_{\mathrm{eq}}$ по-прежнему входит измеренное $\Lambda_T$. Для симплектического отображения оно совпадает с~\eqref{eq:map-xi}, и условия $\xi_x=\xi_y=0$ вместе с $\eta_1=0$ обращают сдвиг в нуль. Жёсткий цилиндр этой последней фразы сам не доказывает.

Надёжная численная реализация суммы есть нормальная форма того же отображения или многооборотное среднее одной орбиты. Оба считают инвариантную кривую, а не эллипс линейной оптики. На игрушке многооборотное среднее сходится к~\eqref{eq:map-id}; одной тысячей оборотов биение тюна ещё не отделено от предела.

\subsection{Почему текущая подстановка не проверка}
\label{sec:map-not}

Задача \texttt{mapping.fox} подставляет в формулу статьи об отображении числа, которые из отображения периода не следуют. Четыре расхождения независимы. Формула статьи есть~\eqref{eq:beta-map}: скобка совпадает со скобкой~\eqref{eq:beta-deq}, предымножитель $1/\eta_0$ имеет обратный знак. Проверка той формулы на отображении должна была бы сравнить среднее $\delta$ орбиты повторения $M$ с правой частью~\eqref{eq:beta-deq}. Подстановка паспортных $\alpha_0$, $\alpha_1$ и $\xi$ этого сравнения не делает.

Первое. Процедура \texttt{MCM} в тексте COSY названа уплотнением $\mathrm{d}L(\mathrm{d}P)/L$. Она переводит каноническое отображение в транспортные координаты, ищет поперечную неподвижную точку и делит приращение пятой транспортной координаты на длину опорной орбиты \texttt{SPOS}. Транспортная шестая координата есть $\delta=\Delta p/p$, не $\delta_K$. Знак транспортной пятой координаты таков, что возвращаемый многочлен есть $-\Delta C/C$. Минус перед коэффициентами в \texttt{mapping.fox} этот знак снимает, и $\alpha_0$, $\alpha_1$ суть коэффициенты пути~\eqref{eq:mag-C} на замкнутой орбите. Извлечённый коэффициент уже является множителем монома, поэтому второй порядок не домножается на $2$. Задачи \texttt{natural.fox} и \texttt{mapping1.fox} домножают на $n!$ и получают производную: их $\alpha_1$ вдвое больше сеничевского коэффициента при $\delta^2$. Паспортное $\eta_1$ собрано без этого множителя. Само по себе $\alpha_1$ есть путь. Высокочастотное поле видит период. Переход от пути к периоду в магните есть~\eqref{eq:mag-eta1}; его и пишет строка \texttt{ETA1} в \texttt{mapping.fox}. Это не чтение $l$.

Второе. Та же магнитная строка применяется к электростатическому и к комбинированному кольцу. Для них $\eta_1$ периода есть~\eqref{eq:es-eta1}, и магнитная комбинация $\alpha_0$, $\alpha_1$ его не даёт. В паспорте электростатического кольца $\eta_1=1{,}07$ есть эта комбинация. Коэффициентом~\eqref{eq:map-eta1} она не является.

Третье. Хроматичность в \texttt{mapping.fox} есть
\begin{equation}
\label{eq:map-chrom}
\xi
=
\bigl(\texttt{MU\_TP}(1)|(0\&0\&0\&0\&1)\bigr)
\left(1+\frac{1}{\gamma_s}\right).
\end{equation}
Процедура \texttt{TP} дифференцирует тюн по параметру, а процедура \texttt{SET\_FOR\_DEUTERONS\_CHROM} заводит этот параметр множителем кинетической энергии опорной частицы. Параметр поэтому есть $\delta_K$. Производная $\mathrm{d}\delta_K/\mathrm{d}\delta_p=(\gamma_s+1)/\gamma_s=1+1/\gamma_s$ переводит $\partial\nu/\partial\delta_K$ в $\partial\nu/\partial\delta_p$. Множитель верен. В электростатике остаётся вопрос, в какой точке потенциала задана опорная кинетическая энергия. Если эта точка лежит при $\Phi=0$, параметр совпадает с $\delta$ главы~\ref{sec:electrostatic}, а потенциал дефлектора уже сидит внутри отображения тюна. Если опорная энергия задана при ненулевом потенциале, тот же множитель переводит локальную кинетическую энергию и перестаёт быть производной по импульсу на зазоре.

Четвёртое относится к спину и разобрано в следующем разделе: коэффициенты \texttt{DAPEE} не являются ни $\Lambda_T$, ни проверкой того, что при $\xi=\eta_1=0$ спиновый тюн перестаёт зависеть от амплитуды.

\subsection{Спиновый тюн}
\label{sec:map-spin}

В магнитном кольце спиновый тюн вдоль замкнутой орбиты есть $G\gamma$. Сдвиг равновесной энергии меняет его на
\begin{equation}
\label{eq:map-spin-eq}
\Delta\nu_s
=
G\beta_s^2\gamma_s\,\delta_{\mathrm{eq}}
+\Delta\nu_s^{\mathrm{direct}},
\end{equation}
где $\mathrm{d}\gamma/\mathrm{d}\delta=\beta_s^2\gamma_s$ и $\delta_{\mathrm{eq}}$ берётся из~\eqref{eq:beta-deq}. Первый член есть канал удлинения орбиты. Второй есть набег уравнения Томаса--Баргмана--Мишеля--Телегди по бетатронной траектории: частица в квадруполе видит радиальные и вертикальные поля, которых нет на оси, и этот набег не сводится к $G\gamma$ опорной орбиты.

Процедура \texttt{TSS} при ключе $0$ возвращает спиновый тюн многочленом по исходным орбитальным координатам, без переноса в неподвижную точку. \texttt{DAPEE} с индексами $11$, $33$ и $66$ извлекает множители $x^2$, $y^2$ и $\delta_K^2$ этого многочлена. Это не производные по действиям. В \texttt{mapping.fox} решётка считается при включённом зазоре. Тогда $\delta_K$ не сохраняется, продольным инвариантом является действие сепаратрисы, и сдвиг центра по $\delta$ уже сидит в амплитудном тюне нормальной формы. Коэффициент при $x^2$ содержит и $G\beta_s^2\gamma_s\,\partial\delta_{\mathrm{eq}}/\partial J_x$, и прямой член.

Обращение $\xi$ и $\eta_1$ в нуль обращает в нуль первый член~\eqref{eq:map-spin-eq}, если $\Lambda_T$ связано с $\xi$ тождеством~\eqref{eq:map-xi}. Прямой член при этом остаётся. Статья 2014 года сводит декогеренцию спина к сдвигу средней энергии: мультиполь меняет средний радиус, длина орбиты меняет равновесную энергию, энергия меняет $G\gamma$~\cite{Senichev2014}. Прямого вклада радиального поля квадруполя в том выводе нет. Формула~\eqref{eq:beta-map} есть тот же энергетический канал, с обратным знаком предымножителя. Численная карта спина в статье об отображении есть полный многочлен \texttt{TSS}, поэтому её нули не обязаны лежать при $\xi_x=\xi_y=\eta_1=0$. Статья это и показывает на всех трёх кольцах. Отдельный вопрос, который здесь не решается: какой вес в $\Delta\nu_s^{\mathrm{direct}}$ имеет сдвиг центра~\eqref{eq:map-centre} относительно члена, квадратичного по бетатронной координате уже на оси квадруполя. Гашение пути эти два вклада не уравнивает, потому что спиновый интеграл поля и интеграл $x/\rho$ суть разные веса одной траектории.

\subsection{Протокол и три кольца}
\label{sec:map-rings}

Проверка~\eqref{eq:beta-deq} и~\eqref{eq:es-deq} на отображении COSY для \texttt{magnetic\_2}, \texttt{electrostatic\_mh} и \texttt{Nuclotron\_16} состоит из двух проходов.

Статическое отображение считается при выключенном зазоре, \texttt{OV 3 3 0}. Длина $C_s$ есть \texttt{SPOS} после решётки. Время на замкнутой орбите даёт \texttt{ECM}; его мономы $e_1$, $e_2$ подставляются в~\eqref{eq:map-eta0} и~\eqref{eq:map-eta1}. Для контроля рядом печатаются $(5|6)$ и $(5|66)$ в начале координат: расхождение есть~\eqref{eq:map-e1} и~\eqref{eq:map-e2}. Путь даёт \texttt{MCM}. Минус перед мономами по транспортному $\delta$ есть $\alpha_0$ и $\alpha_1$; магнитная комбинация~\eqref{eq:mag-eta1} печатается отдельно и с $\eta_1$ из времени сравнивается, а не подменяет его. Твисс берётся из поперечных блоков $\mathrm{ME}$. Квадратичная форма~\eqref{eq:map-ell} и центр~\eqref{eq:map-centre}, включая горизонтальный центр от вертикального эллипса, дают $\Lambda_T/J_x$ и $\Lambda_T/J_y$. Хроматичность считается вторым проходом, \texttt{OV 3 2 1}, и переводится множителем~\eqref{eq:map-chrom}. Рядом печатается правая часть~\eqref{eq:map-xi}. Тождество требует совпадения двух колонок в порядке разложения; запуск ниже показывает, что для $\eta$ это так, а для $\Lambda_T$ на \texttt{magnetic\_2} нет.

Эталон равновесия есть другой запуск, с зазором. Орбита повторения $M$ усредняется по оборотам, и среднее $\delta$ сравнивается с~\eqref{eq:es-deq}. Однооборотные коэффициенты этого среднего не заменяют: они проверяют $\eta$ и $\Lambda_T$, а не захват. Задача \texttt{equilibrium\_check.fox} пишет коэффициенты и обе колонки $\Lambda_T$. Многооборотный эталон в неё не включён. Модуль комбинированного кольца в файле решётки сохранён под именем \texttt{seq\_Nuclotron\_16}.

Задача выполнена для \texttt{magnetic\_2} при выключенных секступолях. Длина $C_s=117{,}388\,\mathrm{м}$. Уплотнение пути $\alpha_0=0{,}23036$, $\alpha_1=0{,}33953$. Время на замкнутой орбите даёт $\eta_0=-0{,}53385$ и $\eta_1=1{,}01779$. Магнитная формула~\eqref{eq:mag-eta1} на тех же $\alpha_0$ и $\alpha_1$ даёт то же $\eta_1$ во всех напечатанных знаках: для чистого магнита путь и время согласованы пересчётом~\eqref{eq:map-eta0}, \eqref{eq:map-eta1}. Паспортные $0{,}230$ и $1{,}02$ суть эти числа, округлённые. В точке старта дисперсия мала, и $(5|6)$ отличается от $e_1$ в пятом знаке. Хроматичность $\xi_x=-1{,}841$, $\xi_y=-1{,}848$ совпадает с паспортом. Квадратичная форма времени, центр при нулевом секступоле не меняет её и в седьмом знаке, даёт $\Lambda_T/J_x=7{,}882\times 10^{-2}\,\mathrm{м}^{-1}$ и $\Lambda_T/J_y=0{,}1194\,\mathrm{м}^{-1}$. Правая часть~\eqref{eq:map-xi} есть $9{,}853\times 10^{-2}$ и $9{,}890\times 10^{-2}\,\mathrm{м}^{-1}$. Отношения равны $0{,}800$ и $1{,}21$. Тождество на этом отображении третьего порядка плоскость за плоскостью не сходится, хотя на игрушечных гамильтоновых киках сходится. Расхождение остаётся открытым: оно не сводится к пропущенному центру и не является общим множителем двух плоскостей. В таблицу ниже магнитное $\eta$ берётся из этого запуска; бетатронная колонка тождества в равновесие не подставляется, пока расхождение не снято.

Паспортные числа трёх колец собраны в таблице~\ref{tab:map-rings}. Магнитная колонка $\eta$ есть~\eqref{eq:mag-eta0} и~\eqref{eq:mag-eta1} на паспортных $\alpha_0$ и $\eta_{1,\mathrm{nat}}$. Правильная колонка есть период. Для чистого магнита они совпадают. Для двух других колец правильное $\eta_1$ из $l$ не извлекалось: в сохранённых выдачах нет мономов \texttt{ECM}. Знак $\delta_{\mathrm{eq}}$ указан при $\Lambda_T>0$ и $\eta_1>0$ по формуле этой заметки. Знак статьи есть знак предымножителя~\eqref{eq:beta-map} на той же скобке: он противоположен. Подстановка одних и тех же $\eta$ и $\Lambda_T$ даёт значения, равные по модулю; расхождение двух формул равно удвоенному истинному сдвигу. На числах главы~\ref{sec:ch2-num} при $\delta_m=10^{-3}$ истинный сдвиг есть $1{,}140\times 10^{-6}$, формула статьи даёт $-1{,}140\times 10^{-6}$.

\begin{table}[t]
\centering
\caption{Линейный и квадратичный коэффициенты периода трёх колец. Для магнита правильная колонка есть запуск \texttt{equilibrium\_check.fox}; паспортное $\eta_1=1{,}02$ есть её округление. Правильное $\eta_1$ электростатики и комбинированного кольца из времени не извлекалось.}
\label{tab:map-rings}
\begin{tabular}{lcccccc}
\hline
Кольцо & $\eta_0^{\mathrm{mag}}$ & $\eta_0$ & $\eta_1^{\mathrm{mag}}$ & $\eta_1$ & знак $\delta_{\mathrm{eq}}$ & знак статьи \\
\hline
\texttt{magnetic\_2} & $-0{,}534$ & $-0{,}5338$ & $1{,}02$ & $1{,}0178$ & $+$ & $-$ \\
\texttt{electrostatic\_mh} & $-0{,}545$ & $-0{,}483$ & $1{,}07$ & нет & $+$ & $-$ \\
\texttt{Nuclotron\_16} & $-0{,}758$ & $-0{,}758$ & $0{,}927$ & нет & $+$ & $-$ \\
\hline
\end{tabular}
\end{table}

Дейтрон \texttt{magnetic\_2}, $\gamma_s=1{,}143914$, $\alpha_0=0{,}230$. Оба $\eta_0$ суть $\alpha_0-1/\gamma_s^2=-0{,}534211$. Паспортное $\eta_1=1{,}02$ для магнита и есть~\eqref{eq:mag-eta1}.

Протон \texttt{electrostatic\_mh}, $\gamma_s=1{,}248329$, $\alpha_0=0{,}0964$. Магнитная формула даёт $\eta_0=-0{,}54531$, электростатическая~\eqref{eq:es-eta0} даёт $\eta_0=-0{,}48345$, как в~\eqref{eq:es-mh}. Оба отрицательны. Паспортное $\eta_1=1{,}07$ остаётся магнитной комбинацией.

Комбинированное кольцо, тот же дейтрон, $\alpha_0=0{,}00648$, паспортное $\eta_1=0{,}927$. Магнитный знаменатель $\eta_0=\alpha_0-1/\gamma_s^2=-0{,}757731$ отрицателен и близок к $-1/\gamma_s^2$: кольцо почти изохронно по геометрии дуг. В фильтре Вина опорная частица идёт прямо, электрический и магнитный радиусы процедуры \texttt{WIEN} равны, но частица с $x\neq 0$ имеет потенциал. Линейный коэффициент периода поэтому есть~\eqref{eq:es-eta0-gen} с $\alpha_0^{E}\neq 0$ на фильтрах, хотя геометрическая кривизна опорной орбиты там равна нулю. Четыре фильтра на период и восемь периодов дают $52{,}3\,\mathrm{м}$ при длине $261{,}8\,\mathrm{м}$. При поле $132\,\mathrm{кВ/см}$ электрический радиус жёсткости около $38\,\mathrm{м}$. Поправка $\alpha_0^{E}/\gamma_s^2$ при дисперсии в фильтре порядка нескольких метров есть сотые и знак $\eta_0$ не меняет; это оценка масштаба, не интеграл по таблице Твисса. Квадратичный коэффициент и $\Lambda_T$ внутри фильтра содержат потенциал главы~\ref{sec:electrostatic}: $\eta_1$ получает добавку~\eqref{eq:es-eta1} на длине фильтров, а $\Lambda_T$ расходится с геометрическим удлинением на $\Lambda_\Phi$, потому что скорость при $\delta=0$ зависит от $x$. Дуги остаются магнитными, и на них $\Lambda_T=\Lambda_C$. Паспортное $\eta_1=0{,}927$ этой добавки не видит.

Эмиттанс в~\eqref{eq:beta-map} при $L=C_s$ совпадает с $\Lambda$ главы~\ref{sec:ch2}, если $\varepsilon=2J$. Среднеквадратичный эмиттанс $\varepsilon_{\mathrm{rms}}=\langle J\rangle$, подставленный в тот же множитель $\pi/L$, даёт половину $\Lambda_T$. Среднее по пучку заменяет $J$ на $\langle J\rangle$ в уже собранном $\Lambda_T$, а не в определении $\varepsilon$ формулы статьи.
